\documentclass[12pt]{article}
\usepackage[utf8]{inputenc}
\usepackage[T1]{fontenc}
\usepackage{hyperref}
\title{Southeastern Tepehuan Language Resources for K-12 Curriculum Development}
\author{Compiled by Claude Sonnet 5.5 (claude-sonnet-5-5) for Andr\'as Kornai}
\date{}
\begin{document}
\maketitle

\section{Introduction}
Southeastern Tepehuan (ISO 639-3 stp; Glottocode sout2976; Uto-Aztecan) is spoken in Durango, Mexico. It is also called O'dam. The best-documented variety is that of Santa Mar\'ia Ocot\'an, in the municipality of El Mezquital. This report is based on a small number of web searches, and little material was found. Nothing found here is a ready-made K-12 curriculum. One source (a language snapshot, see below) says the language is taught in rural schools. This differs from the survey's reference description, and it was not verified independently.

\section{Primary Educational Resources}
No primers, textbooks or school curricula were found in the searches. The dictionary listed below is the closest thing to a teaching-oriented resource.

\section{Linguistic and Reference Materials}
\begin{itemize}
\item Tepehuan Dictionary of Santa Mar\'ia Ocot\'an, Durango (SIL/ILV, ed.\ Willett and Willett, 2015), online on Webonary. It has a Tepehuan--Spanish section and more than five thousand words. The online edition is described as a work in progress. \url{https://www.webonary.work/tepehuano-sureste/overview/introduction/?lang=en}
\item The dictionary's grammar section, covering the phonology and grammar of Southeastern Tepehuan. \url{https://www.webonary.work/tepehuano-sureste/language/grammar/?lang=en}
\item Dictionary guide, ``entries explained''. \url{https://www.webonary.work/tepehuano-sureste/overview/entries-explained/?lang=en}
\item Dictionary foreword. \url{https://www.webonary.work/tepehuano-sureste/overview/foreword/}
\item Garc\'ia Salido and Everdell (2020), language snapshot in \emph{Language Documentation and Description} 19. It describes the documentation state of Tepehuan varieties and notes that the reference grammar is limited. \url{https://lddjournal.org/article/1243/galley/2486/download}
\item General overview (Wikipedia). \url{https://en.wikipedia.org/wiki/Tepehu%C3%A1n_language}
\end{itemize}

\section{Religious and Community Resources}
No scripture translations or community organisations were identified in the searches.

\section{Digital and Audio Resources}
The Webonary dictionary is the only digital resource confirmed. The Endangered Languages Project has a record for the language. \url{https://endangeredlanguages.com/node/112162}

\section{Curriculum Development Context}
Documentation is mostly lexical. A curriculum would have to build on the dictionary and on new grammar and materials work, and spelling conventions vary between sources. The searches found no information on teachers, school programmes or official policy.

\end{document}